\section*{Bab \ref{ch:b3:structural-biology} --- Biologi Struktur Protein}

\begin{solution}{exo:b3:structural-biology:1}
Primer: urutan $141$ residu sebuah rantai $\alpha$.
Sekunder: kedelapan heliks $\alpha$-nya. Tersier: \omterm{def:b3:structural-biology:levels}{lipatan} globin satu
rantai yang membungkus hemnya. Kuaterner: tetramer $\alpha_{2}\beta_{2}$
dan cara dimernya berputar satu sama lain.
\end{solution}

\begin{solution}{exo:b3:structural-biology:2}
$36\times \qty{0.15}{nm} = \qty{5.4}{nm}$; $36/3.6 = 10$ putaran;
$36 - 4 = 32$ ikatan hidrogen.
\end{solution}

\begin{solution}{exo:b3:structural-biology:3}
Bahwa struktur asalnya, termasuk pasangan yang benar dari empat
disulfida di antara $105$ kemungkinannya, dipulihkan secara spontan dari
rantai yang tak terlipat: urutannya memuat seluruh keterangan yang
dibutuhkan, dan keadaan asalnya adalah minimum termodinamikanya. Mengerjakannya
di dalam tabung reaksi dengan protein murni menyingkirkan cetakan, enzim, atau
permesinan sel mana pun sebagai sumber keterangannya.
\end{solution}

\begin{solution}{exo:b3:structural-biology:4}
NMR: protein kecil, di bawah kira-kira \qty{40}{kDa}, di dalam larutan.
Kristalografi: ukuran berapa pun yang dapat mengkristal, dari peptida sampai
ribosom. \omterm{def:b3:structural-biology:methods}{Mikroskopi elektron krio}: kompleks besar, di atas kira-kira
\qty{50}{kDa}, tanpa kristal. NMR melaporkan gerak secara langsung (dan
krio-EM dapat memilah molekulnya menjadi golongan konformasi); sebuah struktur
kristal adalah rata-rata yang statis.
\end{solution}

\begin{solution}{exo:b3:structural-biology:5}
$2^{150} \approx 1.4\times 10^{45}$ konformasi, $1.4\times 10^{33}$
s $\approx 5\times 10^{25}$ tahun. $3^{20} \approx 3.5\times 10^{9}$,
$3.5$ ms: peptidanya dapat menelusuri seluruhnya dalam satu detik;
proteinnya tidak dapat, bahkan sepanjang umur alam semesta.
\end{solution}

\begin{solution}{exo:b3:structural-biology:6}
$Y = [\alpha(1+\alpha) + Lc\alpha(1+c\alpha)]/[(1+\alpha)^{2} +
L(1+c\alpha)^{2}]$. $\alpha = 1$: $3.01/106 = 0.028$ (satu tapak
$0.50$). $\alpha = 10$: $121/242 = 0.50$ ($0.91$). $\alpha = 100$:
$\num{10300}/\num{10601} = 0.97$ ($0.99$). Kurva MWC-nya jauh tertinggal
di belakang satu tapak pada ligan rendah, lalu naik curam --- dari $0.03$
menjadi $0.97$ sepanjang dua dasawarsa ketika satu tapaknya bergerak dari
$0.5$ menjadi $0.99$: ketajaman itulah \omterm{def:b3:structural-biology:allostery}{kekooperatifannya}.
\end{solution}

\begin{solution}{exo:b3:structural-biology:7}
Bisfosfogliserat hanya mengikat \omterm{def:b3:structural-biology:allostery}{keadaan T} (di dalam rongga pusatnya), sehingga
ia memantapkan T: di dalam modelnya ia menaikkan $L$, yakni nisbah $T/R$
protein yang tak berligan, sehingga menggeser kurvanya ke kanan --- kegemaran
lebih rendah, lebih banyak oksigen dilepaskan pada tekanan jaringan. Sel darah
merah yang disimpan kehilangan senyawa itu, $L$-nya turun, kurvanya bergeser
ke kiri, dan darah transfusi menahan oksigennya alih-alih melepaskannya di
dalam jaringan, sampai selnya mensintesisnya kembali sepanjang sehari.
\end{solution}

\begin{solution}{exo:b3:structural-biology:8}
$d = \lambda/(2\sin\theta) = 0.1/(2\times 0.259) = \qty{0.19}{nm}$: tulang
punggung dan setiap rantai sampingnya tertempatkan, atom demi atom terpisahkan.
Bagi \qty{0.12}{nm}: $\sin\theta = 0.1/0.24 = 0.417$, $\theta =
\qty{25}{\degree}$.
\end{solution}

\begin{solution}{exo:b3:structural-biology:9}
Sebuah lisina bermuatan yang terkubur di antara rantai samping nonpolar
memakan puluhan kilojoule tiap mol (muatannya tak tersolvasi, penguburan
hidrofob leusinanya hilang), lebih besar daripada seluruh $\Delta G$
pelipatannya: proteinnya membuka \omterm{def:b3:structural-biology:levels}{lipatan} atau intinya menata ulang. Pada
permukaannya, lisinanya terhidrasi sebagaimana pada keadaan tak terlipat dan
tak ada yang hilang. $\qty{-30}{kJ/mol} \approx -12RT$ berarti $K = e^{12} \approx 1.6\times
10^{5}$: satu molekul dari \num{160000} tak terlipat pada tiap saat, dan
seluruh marginnya setara dua atau tiga ikatan hidrogen dari ratusan ikatan
--- senilai satu mutasi saja.
\end{solution}

\begin{solution}{exo:b3:structural-biology:10}
Spesi $R$-nya berjumlah $[R_{0}](1+\alpha)^{n}$ dan spesi $T$-nya
$[R_{0}]L(1+c\alpha)^{n}$, sehingga $\bar R = (1+\alpha)^{n}/[(1+\alpha)^{n}
+ L(1+c\alpha)^{n}]$. Untuk $c \to 0$ suku $T$-nya menjadi $L$, dan $\bar R =
1/2$ ketika $(1+\alpha)^{n} = L$, yakni $\alpha = L^{1/n} - 1$. Bagi
hemoglobin $(3\times 10^{5})^{1/4} = 23.4$, sehingga $\alpha \approx 22$:
\qty{22}{mmHg}, dekat dengan $P_{50}$-nya \qty{26}{mmHg} --- sakelar
konformasinya dan setengah-jenuhnya nyaris berimpit.
\end{solution}

\begin{solution}{exo:b3:structural-biology:11}
Hipotesis itu ditentang karena setiap agen penular yang dikenal bereplikasi
lewat asam nukleat, dan sebuah protein tak dapat menyalin urutannya. Bila
sebuah asam nukleat ternyata dibutuhkan, atau protein murni terbukti tak dapat
menularkannya, hipotesisnya terbantah. Dukungannya: daya infeksinya bertahan
terhadap nuklease, cahaya ultraungu, dan radiasi pengion pada dosis yang
memusnahkan genom mana pun, tetapi musnah oleh zat pendenaturasi protein dan
protease; PrP rekombinan yang dilipat di luar sel menjadi fibril menularkan
penyakitnya kepada hewan, dengan sifat galur yang sama pada penyalurannya.
Sebuah benih dibutuhkan karena perubahan sebuah monomer normal terlalu lambat
bila sendirian --- inti struktur silang-$\beta$-nya harus terbentuk lebih
dahulu, dan itulah peristiwa yang membatasi lajunya --- sedangkan ujung fibril
yang telah terbentuk mengubah monomernya dengan cepat: bentuknya disalin lewat
pertumbuhan tercetak.
\end{solution}

\begin{solution}{exo:b3:structural-biology:12}
Banyak urutan melipat menjadi struktur yang sama: yang dituntut sebuah \omterm{def:b3:structural-biology:levels}{lipatan}
adalah pola kedudukan hidrofob dan polar, beberapa glisina dan prolina
pada belokannya, serta residu yang mengerjakan tugasnya; segala yang lain pada
permukaannya dapat berubah tanpa denda, sehingga mutasi menumpuk di situ
sementara seleksi mempertahankan \omterm{def:b3:structural-biology:levels}{lipatan} dan fungsinya. Karena itu residu yang
lestari adalah intinya (lebih sebagai pola daripada sebagai jati diri), residu
katalitik dan pengikatnya, serta glisina, prolina, dan sisteina strukturalnya.
Metode \omterm{def:b3:bioinformatics:msa}{profil} membobot justru lajur itu sehingga mengenali anggota famili pada
keidentikan \qty{15}{\%}; perancangan protein menjalankan logikanya terbalik,
dengan memilih pola hidrofob bagi \omterm{def:b3:structural-biology:levels}{lipatan} sasarannya lalu membiarkan
permukaannya menjadi apa saja.
\end{solution}

\begin{solution}{pb:b3:structural-biology:1}
\textbf{1.} $0.75\times 150 \approx 112$ residu heliks; $112\times
\qty{0.15}{nm} = \qty{17}{nm}$; $112/3.6 \approx 31$ putaran.
\textbf{2.} $112 - 8\times 4 = 80$ ikatan hidrogen.
\textbf{3.} Massanya $150\times 110 = \qty{16500}{Da} = 2.7\times 10^{-20}$
g; volumenya $2.7\times 10^{-20}\times 0.73 = 2.0\times 10^{-20}$ cm$^{3}$
$= \qty{20}{nm^3}$; jari-jarinya $(3V/4\pi)^{1/3} = \qty{1.7}{nm}$.
\textbf{4.} Terjulur $150\times 0.36 = \qty{54}{nm}$ terhadap garis tengah
\qty{3.4}{nm}: faktor $16$.
\textbf{5.} Residu heliksnya \qty{75}{\%} panjang terjulurnya,
$112\times 0.36 = \qty{40}{nm}$; heliksnya menekannya menjadi
\qty{17}{nm}, faktor $0.36/0.15 = 2.4$.
\textbf{6.} $40\times 4 = \qty{160}{kJ/mol}$ diperoleh, $150\times 0.9 =
\qty{135}{kJ/mol}$ hilang: neto $\Delta G \approx \qty{-25}{kJ/mol}$, di
dalam rentang tipisnya.
\textbf{7.} $3^{150} \approx 4\times 10^{71}$; $4\times 10^{59}$ s
$\approx 10^{52}$ tahun.
\textbf{8.} $10^{-2}\times 10^{12} = 10^{10}$ konformasi, sebesar bagian
$3\times 10^{-62}$ dari cacahannya.
\textbf{9.} Heliksnya dalam \qty{1}{\micro s} (secara paralel); $8! = \num{40320}$
pengemasan pada $10^{6}$ tiap detik memakan \qty{40}{ms}: seluruhnya sekitar
\qty{40}{ms}, yakni orde besar yang teramati. Ya: memecahkan bagiannya
secara bebas lalu keseluruhannya adalah sebuah corong --- ruang telusurnya
difaktorkan, bukan dihabiskan.
\textbf{10.} Harapan cacah daurnya $1/0.3 = 3.3$, kira-kira \qty{33}{s};
sesudah enam daur tersisa $0.7^{6} = 0.12$ yang tak terlipat.
\textbf{11.} $\Delta G = -RT\ln K$ dengan $K = 10^{4} \to 10^{2}$: dari
$-23.7$ menjadi \qty{-11.9}{kJ/mol}, yakni kehilangan \qty{12}{kJ/mol}
kemantapan ($RT = \qty{2.58}{kJ/mol}$, $\ln 100 = 4.6$).
\textbf{12.} Penggumpalan dimulai dengan sebuah inti, yang laju
pembentukannya naik sebagai pangkat tinggi kepekatan spesi tak terlipat yang
rawan menggumpal; seratus kali lebih banyak spesi itu menaikkan peluang
terbentuknya sebuah benih sepanjang satu umur hidup sampai beberapa orde besar,
dan begitu sebuah benih ada, pertumbuhannya cepat dan mencetak dirinya sendiri.
\textbf{13.} $\alpha = 100$: $Y = 1.054\times 10^{8}/1.089\times 10^{8} =
0.97$. $\alpha = 20$: pembilangnya $185\,220 + 103\,680 = 288\,900$,
penyebutnya $194\,481 + 622\,080 = 816\,561$: $Y = 0.35$.
\textbf{14.} $\bar R = 1.041\times 10^{8}/1.089\times 10^{8} = 0.96$ di
dalam paru-paru; $194\,481/816\,561 = 0.24$ di dalam ototnya.
\textbf{15.} $0.97 - 0.35 = 0.62$ kemampuannya dilepaskan (dua pertiga
muatannya). Satu tapak: $100/126 - 20/46 = 0.79 - 0.43 = 0.36$.
\textbf{16.} Kemampuannya $150\times 1.34 = \qty{200}{mL/L}$; ke ototnya
$0.62\times 200 \approx \qty{125}{mL/L}$. Saat istirahat $(0.97 - 0.78)\times
200 = \qty{38}{mL/L}$; $250/38 \approx \qty{6.5}{L/min}$ aliran
darah --- seorde dengan curah jantung saat istirahat.
\textbf{17.} $L = 3\times 10^{6}$, $\alpha = 40$: pembilangnya $2.76\times
10^{6} + 3.29\times 10^{6} = 6.05\times 10^{6}$, penyebutnya $2.83\times
10^{6} + 11.5\times 10^{6} = 14.4\times 10^{6}$: $Y = 0.42$ alih-alih
$0.78$ --- jauh lebih banyak oksigen dilepaskan saat istirahat, yakni
penyesuaian darah terhadap ketinggian dan otot terhadap kerja.
\textbf{18.} Janin, $L = 3\times 10^{4}$, $\alpha = 30$: pembilangnya
$893\,730 + 19\,770 = 913\,500$, penyebutnya $923\,520 + 85\,680 =
1\,009\,200$: $Y = 0.91$. Ibunya pada tekanan yang sama: pembilangnya
$893\,730 + 197\,730 = 1\,091\,460$, penyebutnya $923\,520 + 856\,830
= 1\,780\,350$: $Y = 0.61$. Pada tekanan yang sama darah janin menahan lebih
banyak oksigen, sehingga oksigennya mengalir dari ibu ke janin melintasi plasenta.
\textbf{19.} $(10^{5}\,\text{nm}/6\,\text{nm})^{3} = 4.6\times 10^{12}$
sel; $1.9\times 10^{13}$ molekul.
\textbf{20.} $d = 0.1/(2\sin 14.5^{\circ}) = \qty{0.20}{nm}$: ya,
setiap rantai samping dan atomnya tertempatkan.
\textbf{21.} $\sin\theta = 0.1/0.3$, $\theta = \qty{19.5}{\degree}$. Di
dalam kristal yang tak teratur, molekulnya tidak sebaris persis, gelombang
yang terhambur pada sudut tinggi --- yang menyandikan rincian halus --- tak
lagi berpadu sefase, dan bintiknya memudar pada sudut tinggi; hanya bintik
berdaya pisah rendah yang bertahan.
\textbf{22.} Memparuhkan $d$ menuntut isyarat dua kali lipat, sehingga empat
kali partikelnya: sekitar \num{40000} (lebih banyak dalam praktiknya, karena
galat lain ikut masuk).
\textbf{23.} Protein membran, yang jarang mengkristal dari larutan deterjen
dan dapat dicitra di dalam nanocakram lipid; kompleks besar yang lentur ---
ribosom, spliseosom, saluran ion beserta pengaturnya --- yang
keheterogenannya mencegah pengkristalan tetapi partikelnya dapat dipilah
menjadi golongan konformasi.
\textbf{24.} Hem dan oksigen yang terikat beserta geometrinya yang persis,
air dan ion yang tertata, konformasi keduanya (T terhadap R), serta bukti
bahwa \omterm{def:b3:structural-biology:levels}{lipatannya} benar --- sebuah ramalan hanyalah satu model tanpa ligan
tanpa pemeriksaan percobaan.
\textbf{25.} Pelepasan $0.62$ kemampuannya (satu tapak $0.36$); waktu
Levinthalnya sekitar $10^{52}$ tahun; sebuah peta pada \qty{0.20}{nm}.
\end{solution}
